% Author: Fabio Aurelio D'Asaro <fabioaurelio.dasaro at univr.it>
% License: CC BY-NC 4.0 (https://creativecommons.org/licenses/by-nc/4.0/).
% University logos and marks are excluded; no rights over them are granted.
% Beavr 2.0. Built on LaTeX Beamer and PGF/TikZ.

\documentclass{beamer} % Add [aspectratio=169] for widescreen.
\usepackage{iftex}
\ifPDFTeX
  \usepackage[T1]{fontenc}
  \usepackage[utf8]{inputenc}
\fi

% 0: orange/red; 1: gray/blue; 2: light blue/ochre; 3: gray.
\def\beavrcolor{0}
% \def\beavrlogo{logo_univr.png} % Set to {} for a logo-free cover.
% \def\beavrlogowidthtitle{4.6cm}
% \def\beavrfooterlogo{} % Omit the footer logo if desired.
\usetheme{beavr}
\usetikzlibrary{arrows.meta}

% Optional short labels appear in the footer.
\title[Introducing Beavr 2.0]{Introducing Beavr 2.0: a Beamer template for UniVR}
\subtitle{\LaTeX{} presentations}
\author[J. Doe]{John Doe}
% \institute[UniVR]{Universit\`a di Verona}
\beavremail{john.doe@univr.it}
\date{2 October 2026} % Replace with your event date, or use \today.
\hypersetup{pdfauthor={John Doe},pdftitle={Introducing Beavr 2.0: a Beamer template for UniVR}}

\begin{document}

% The cover is excluded from the frame count.
\begin{frame}[plain,noframenumbering]
  \titlepage
\end{frame}

\begin{frame}{There is no largest prime number}
  \framesubtitle{A proof by \textit{reductio ad absurdum}}
  \begin{theorem}
    There is no largest prime number.
  \end{theorem}
  \begin{proof}
    \begin{enumerate}
      \item Suppose $p$ is the largest prime number.
      \item Let $q$ be the product of all primes up to $p$.
      \item No prime up to $p$ divides $q+1$.
      \item Since $q+1>1$, it has a prime divisor greater than $p$:
        a contradiction.
    \end{enumerate}
  \end{proof}
\end{frame}

\begin{frame}{Lists, examples and alerts}
  \begin{itemize}
    \item A first-level item
      \begin{itemize}
        \item A first subitem
        \item A second subitem
      \end{itemize}
    \item Back to the first level
  \end{itemize}
  \begin{example}[Scenario 1]
    An example illustrates the main idea.
  \end{example}
  \begin{alertblock}{A point to remember}
    This is an alert!
  \end{alertblock}
\end{frame}

% TikZ example.
\begin{frame}[t]{Model fitting}
  \framesubtitle{An iterative modelling workflow}
  \begin{columns}[T,onlytextwidth]
    \begin{column}{.53\textwidth}
      \begin{block}{The workflow}
        \centering
        \begin{tikzpicture}[
          stage/.style={draw=fgcolor,thick,rounded corners=3pt,
            fill=fgcolor!6!white,text width=1.85cm,
            minimum height=.75cm,align=center,font=\small},
          flow/.style={-{Latex[length=2mm]},thick,draw=fgcolor},
          label/.style={font=\scriptsize,text=fgcolor,fill=none,inner sep=2pt}]
          \node[stage] (data) at (0,0) {Observed data};
          \node[stage,fill=bgcolor] (model) at (2.9,0) {Model $\theta$};
          \node[stage] (predict) at (2.9,-1.35) {Predictions};
          \node[stage] (check) at (0,-1.35) {Evaluation};
          \draw[flow] (data) -- node[label,above] {fit} (model);
          \draw[flow] (model) -- node[label,right] {predict} (predict);
          \draw[flow] (predict.south) -- (2.9,-1.95) --
            node[label,below] {compare} (0,-1.95) -- (check.south);
          \draw[flow] (check.north east) --
            node[label,sloped,above] {update} (model.south west);
        \end{tikzpicture}
      \end{block}
    \end{column}
    \begin{column}{.43\textwidth}
      \begin{itemize}
        \item Start with observations and an initial model.
        \item Compare predictions with the data.
        \item Use the error to refine the model and repeat.
      \end{itemize}
    \end{column}
  \end{columns}
  \vfill
  \begin{exampleblock}{Gradient descent}
    {\centering
      $\displaystyle\theta_{t+1}=\theta_t-\eta\,\nabla_{\theta}\mathcal{L}(\theta_t)$\par}
    \smallskip
    The learning rate $\eta$ controls the size of each update.
  \end{exampleblock}
\end{frame}

\end{document}
